
\subsubsection{6.1 Interpreting the Near-Perfect Null}

The combination of 5-gate all-pass and HR = 1.006 (p = 0.9495, 176 displacement events) produces a clean scientific conclusion: community breadth diversity at benchmark introduction year, measured by paper submission count, does not predict plurality benchmark displacement timing in Papers With Code (2015--2023). This is not a marginal null. The LRT statistic is 0.0040 and $\Delta logL$ = 0.0020, placing the result extremely close to the null prediction of no effect.

\subsubsection{6.2 Why Does the Mechanism Not Operate?}

Four explanations are considered, ordered by assessed plausibility:

\textbf{E1 (HIGH): Predictor construct mismatch.} The prior directional signal (HR = 0.871) was for \textit{score velocity} ($\Delta score$\_lag1\_z) --- a fundamentally different construct from submission-count diversity. Score velocity measures how fast SOTA improved; submission count measures how many teams participated. The current null applies to submission-count breadth only and says nothing about whether score velocity would yield a non-null result at 176 events.

\textbf{E2 (HIGH): Construct validity gap.} \texttt{paper\textbackslash \_url} deduplication is paper-level, not institution-level. If a small number of prolific laboratories dominate submissions, high submission count does not imply broad stakeholder investment. The lock-in mechanism requires \textit{institutional} diversity (distinct teams with sunk costs). The null may reflect inadequacy of the proxy rather than absence of the mechanism.

\textbf{E3 (MEDIUM): Mechanism cancellation.} Lock-in (HR < 1) and saturation (HR > 1) may both operate but cancel across the population. The aggregate HR $\approx$ 1.0 would then reflect a mixture of two opposing pathways, disentangleable only with moderator variables (e.g., task difficulty trajectory).

\textbf{E4 (LOW): Residual temporal confound.} Partial $r^2$ = 0.605 strongly argues against this explanation.

The most parsimonious interpretation combines E1 and E2: the null reflects both the use of a different predictor construct from that which showed a directional signal, and a potential mismatch between submission count and the institutional investment that the lock-in mechanism requires.

\subsubsection{6.3 Implications for Future Work}

\textbf{Score trajectory:} Cox regression with $\Delta score$\_lag1\_z on the full 258-row complete-case panel (176 events) would test whether the prior HR = 0.871 (22 events) survives adequate statistical power.

\textbf{Institutional diversity:} Augmenting the h-e1 pipeline with Semantic Scholar author API to compute unique-institution count at introduction would test whether true stakeholder diversity predicts displacement when submission count does not.

\subsubsection{6.4 Limitations}

\textbf{L1 (Construct validity):} Submission count $\neq$ institutional adopter diversity. The null is valid for this operationalization; it does not rule out better-constructed breadth measures.

\textbf{L2 (25\% data reduction):} 87 of 345 rows were dropped (NaN, primarily unmatched benchmarks). The remaining 258 rows yield 176 displacement events; EPV with M1's three covariates is approximately 58.7, which is adequate by conventional standards. However, potential selection bias exists if unmatched benchmarks differ systematically from matched ones; this was not investigated.

\textbf{L3 (PH assumption):} \texttt{check\textbackslash \_assumptions()} was not completed due to a string-encoding error in \texttt{task\textbackslash \_path} (requires integer-encoding before calling). Visual inspection of Schoenfeld residuals showed no obvious PH violation; quantitative confirmation is pending. Given HR = 1.006, the conclusion is robust to moderate PH misspecification.

\textbf{L4 (Scope):} Results are specific to Papers With Code (CC-BY-SA-4.0), 87 Koch et al. 2021 tasks, 2015--2023. Generalization to other registries or time periods is not established.

\textbf{L5 (Penalizer sensitivity):} The L2 penalty (penalizer = 0.1) was selected for convergence, not optimized for effect estimation. In samples of this size, L2 regularization can shrink coefficients toward zero. The null result HR = 1.006 should be interpreted as ''not distinguishable from null under L2 regularization with penalizer = 0.1.'' Sensitivity analysis with penalizer = 0.0 (unpenalized) and penalizer = 0.5 would confirm the null is not a regularization artifact; these runs were not performed.

\textbf{L6 (KM log-rank p-value):} The Kaplan-Meier Q1 vs Q4 comparison (Section 5.3) is supported only by visual inspection. A numeric log-rank p-value was not extracted in the current pipeline run. Given the near-perfect Cox null, the KM conclusion is consistent but lacks quantitative confirmation.

\subsubsection{6.5 On the Value of Principled Nulls}

A null under poor measurement is uninformative. A null under validated measurement (5-gate all-pass) and adequate statistical power (176 displacement events) definitively rules out a predictor class, saving future researchers from re-investigating the same construct. Submission-count community breadth diversity is such a predictor: well-measured by the FAIL FAST protocol, tested at adequate power, and definitively null.

